\chapter{Geometric Properties of Thought}
\label{ch:thought_geometry}

\papernote{Source paper: \emph{On the Thermodynamic Consequences of Categorical Completion: Geometric Methods for Fluid-Based Hybrid Analog Oscillatory Integrated Circuits} --- \texttt{mekaneck/docs/geometrics\_properties\_of\_cognitive\_processes/}}

\begin{chapterabstract}
Individual thoughts are three-dimensional geometric objects: specific oxygen molecular gas configurations completing oscillatory hole circuits. This chapter characterises thought geometry formally, establishes the oscillatory-categorical equivalence for thought formation, and validates the framework through the olfactory paradigm. Olfactory recognition is shown to proceed via inelastic electron tunnelling spectroscopy rather than shape-based (lock-and-key) docking---a result that (i) validates the categorical completion mechanism and (ii) explains why enantiomers of the same molecule can smell different despite having identical shapes.
\end{chapterabstract}

\input{../mekaneck/docs/geometrics_properties_of_cognitive_processes/geometric-properties-of-thought.tex}
